% from ega2-1-fr.tex
\oldpage[II]{142}
condition b) implique que $f(X)=Y$. Or, soit $y'$ un point de $Y$ et soit $y$
le point générique d'une composante irréductible de $Y$ contenant $y'$~; en
vertu de b), il suffit de montrer que $f^{-1}(y)$ n'est pas vide. Mais on sait
que cette propriété est conséquence du fait que $f$ est dominant et
quasi-compact (\hyperref[I.6.6.5-fr]{I, 6.6.5}).

\begin{corollary}[7.2.2]
\label{II.7.2.2-fr}
Soit $f:X\to Y$ un morphisme d'immersion quasi-compact. Pour que l'espace
sous-jacent $X$ soit fermé dans $Y$, il faut et il suffit qu'il contienne toute
spécialisation (dans $Y$) de chacun de ses points.
\end{corollary}

\begin{proposition}[7.2.3]
\label{II.7.2.3-fr}
Soient $Y$ un préschéma (resp. un préschéma localement noethérien), $f:X\to Y$
un morphisme (resp. un morphisme localement de type fini). Les conditions
suivantes sont équivalentes~:
\begin{enumerate}
\item[\rm a)] $f$ est séparé.
\item[\rm b)] Le morphisme diagonal $X\to X\times_YX$ est quasi-compact, et
pour tout $Y$-préschéma de la forme $Y'=\operatorname{Spec}(A)$, où $A$ est un
anneau de valuation (resp. un anneau de valuation discrète), deux
$Y$-morphismes de $Y'$ dans $X$ qui coïncident au point générique de $Y'$ sont
égaux.
\item[\rm c)] Le morphisme diagonal $X\to X\times_YX$ est quasi-compact, et
pour tout $Y$-préschéma de la forme $Y'=\operatorname{Spec}(A)$, où $A$ est un
anneau de valuation (resp. un anneau de valuation discrète), deux
$Y'$-sections de $X'=X_{(Y')}$ qui coïncident au point générique de $Y'$ sont
égales.
\end{enumerate}
\end{proposition}

L'équivalence de b) et de c) résulte de la correspondance biunivoque entre
$Y$-morphismes de $Y'$ dans $X$ et $Y'$-sections de $X'$
(\hyperref[I.3.3.14-fr]{I, 3.3.14}). Si $X$ est séparé sur $Y$, la condition b)
est vérifiée en vertu de (\hyperref[I.7.2.2.1-fr]{I, 7.2.2.1}), puisque $Y'$
est intègre. Il reste à prouver que b) entraîne que le morphisme diagonal
$\Delta:X\to X\times_YX$ est fermé, et il revient au même de montrer qu'il
vérifie le critère de (\hyperref[II.7.2.2-fr]{7.2.2}). Or, soit $z$ un point de
la diagonale $\Delta(X)$, $z'\neq z$ une spécialisation de $z$ dans
$X\times_YX$. Il existe alors (\hyperref[II.7.1.4-fr]{7.1.4}) un anneau de
valuation $A$ et un morphisme $f$ de $Y'=\operatorname{Spec}(A)$ dans
$X\times_YX$, qui applique le point fermé $a$ de $Y'$ sur $z'$ et le point
générique $b$ de $Y'$ sur $z$~; ce morphisme fait de $Y'$ un
$(X\times_YX)$-préschéma, et \emph{a fortiori} un $Y$-préschéma. Si on compose
avec $f$ les deux projections de $X\times_YX$, on obtient deux $Y$-morphismes
$g_1$, $g_2$ de $Y'$ dans $X$, qui par hypothèse coïncident au point $b$~; ils
sont donc égaux à un même morphisme $g$, ce qui signifie
(\hyperref[I.5.3.1-fr]{I, 5.3.1}) que $f$ se factorise en $f=\Delta\circ g$, et
par suite $z'\in\Delta(X)$. Lorsqu'on suppose $Y$ localement noethérien et $f$
localement de type fini, $X\times_YX$ est localement noethérien
(\hyperref[I.6.6.7-fr]{I, 6.6.7})~; on peut alors reprendre le raisonnement
précédent en y supposant que $A$ est un anneau de valuation discrète, en vertu
de (\hyperref[II.7.1.9-fr]{7.1.9}).

\begin{remark}[7.2.4]
\label{II.7.2.4-fr}
\begin{enumerate}
\item[\rm (i)] L'hypothèse que le morphisme $\Delta$ est quasi-compact est
toujours vérifiée lorsque $Y$ est localement noethérien et $f$ localement de
type fini, car $X\times_YX$ est alors localement noethérien
(\hyperref[I.6.6.4-fr]{I, 6.6.4, (i)}). Dans le cas général, elle signifie aussi
que pour tout recouvrement $(U_\alpha)$ de $X$ par des ouverts affines, les
ensembles $U_\alpha\cap U_\beta$ sont \emph{quasi-compacts}.
\item[\rm (ii)] Pour que $f$ soit séparé, il \emph{suffit} que la condition b)
ou la condition c) soit vérifiée pour un anneau de valuation $A$ qui est
\emph{complet} et dont le corps des restes est \emph{algébriquement clos}~;
cela résulte en effet de la démonstration de
(\hyperref[II.7.2.3-fr]{7.2.3}) et de
(\hyperref[II.7.1.4-fr]{7.1.4}).
\end{enumerate}
\end{remark}

